\documentclass[11pt,a4paper]{article}

\usepackage[utf8]{inputenc}
\usepackage[T1]{fontenc}
\usepackage[greek,english]{babel}
\usepackage[margin=1in]{geometry}
\usepackage{graphicx}
\usepackage{amsmath,amssymb}
\usepackage{enumitem}
\usepackage{parskip}
\usepackage{placeins}
\usepackage{booktabs}
\usepackage{algorithm}
\usepackage{algpseudocode}
\usepackage{hyperref}
\usepackage{float}

\usepackage{listings}
\usepackage{xcolor}

\lstset{
  basicstyle=\ttfamily\small,
  keywordstyle=\color{blue},
  commentstyle=\color{gray}\itshape,
  stringstyle=\color{orange},
  numbers=left,
  numberstyle=\tiny\color{gray},
  breaklines=true,
  frame=single,
  showstringspaces=false,
  language=Python
}

\hypersetup{
  colorlinks=true,
  linkcolor=blue,
  urlcolor=cyan
}

\title{
  \textbf{\foreignlanguage{english}{SLP-NLP Lab 2}}\\[0.5em]
  \large Επεξεργασία Φωνής και Φυσικής Γλώσσας\\
  Εαρινό Εξάμηνο 2025-26
}
\author{Μαντζώρος Γεράσιμος, Παπαδημητρίου Νίκη}
\date{Απρίλιος 2026}

\begin{document}

\selectlanguage{greek}

\maketitle
\tableofcontents
\newpage

% ============================================================
%  ΜΕΡΟΣ Α: ΠΡΟΠΑΡΑΣΚΕΥΗ
% ============================================================
\phantomsection
\section*{Μέρος Α΄: Προπαρασκευή}
\addcontentsline{toc}{section}{Μέρος Α΄: Προπαρασκευή}

Για τα ζητούμενα της προπαρασκευής χρησιμοποιήσαμε το \textbf{\textlatin{[Sentence Polarity Dataset / Semeval 2017 Task4-A]}}
% TODO: επιλέξτε ένα dataset και διαγράψτε το άλλο

\phantomsection
\subsection*{Ζητούμενο 1 -- Κωδικοποίηση Επισημειώσεων}
\addcontentsline{toc}{subsection}{Ζητούμενο 1 -- Κωδικοποίηση Επισημειώσεων}

% Εκτυπώστε τα πρώτα 10 labels από τα δεδομένα εκπαίδευσης και τις αντιστοιχίες τους σε αριθμούς.

\begin{lstlisting}[caption={\textlatin{Έξοδος Ζητούμενου 1 (main.py:EX1)}}]
% TODO: Επικολλήστε εδώ την έξοδο της εκτύπωσης των 10 πρώτων labels
\end{lstlisting}

\phantomsection
\subsection*{Ζητούμενο 2 -- Λεκτική Ανάλυση (\textlatin{Tokenization})}
\addcontentsline{toc}{subsection}{Ζητούμενο 2 -- Λεκτική Ανάλυση (Tokenization)}

% Εκτυπώστε τα πρώτα 10 παραδείγματα από τα δεδομένα εκπαίδευσης (dataloading.py:EX2).

\begin{lstlisting}[caption={Έξοδος Ζητούμενου 2 (dataloading.py:EX2)}]
% TODO: Επικολλήστε εδώ την έξοδο των 10 πρώτων tokenized παραδειγμάτων
\end{lstlisting}

\phantomsection
\subsection*{Ζητούμενο 3 -- Κωδικοποίηση Παραδειγμάτων (Λέξεων)}
\addcontentsline{toc}{subsection}{Ζητούμενο 3 -- Κωδικοποίηση Παραδειγμάτων}

% Υλοποίηση __getitem__ της κλάσης SentenceDataset (dataloading.py:EX3)
% Εκτυπώστε 5 παραδείγματα στην αρχική τους μορφή και όπως τα επιστρέφει η κλάση.

\begin{lstlisting}[caption={Έξοδος Ζητούμενου 3 (dataloading.py:EX3)}]
% TODO: Επικολλήστε εδώ τα 5 παραδείγματα όπως τα επιστρέφει η SentenceDataset
\end{lstlisting}

\phantomsection
\subsection*{Ζητούμενο 4 -- \textlatin{Embedding Layer}}
\addcontentsline{toc}{subsection}{Ζητούμενο 4 -- Embedding Layer}

% Συμπληρώστε τα κενά στις θέσεις models.py:EX4
% Απαντήστε στα παρακάτω ερωτήματα:

\paragraph{Γιατί αρχικοποιούμε το \textlatin{embedding layer} με τα προ-εκπαιδευμένα \textlatin{word embeddings};}

% TODO: Η αρχικοποίηση με προ-εκπαιδευμένα embeddings δίνει στο μοντέλο ...

\paragraph{Γιατί κρατάμε παγωμένα τα βάρη του \textlatin{embedding layer} κατά την εκπαίδευση;}

% TODO: Παγώνοντας τα βάρη αποφεύγουμε ...

\phantomsection
\subsection*{Ζητούμενο 5 -- \textlatin{Output Layer(s)}}
\addcontentsline{toc}{subsection}{Ζητούμενο 5 -- Output Layer(s)}

% Συμπληρώστε τα κενά στις θέσεις models.py:EX5
% Απαντήστε στα παρακάτω ερωτήματα:

\paragraph{Γιατί βάζουμε μη γραμμική συνάρτηση ενεργοποίησης στο προτελευταίο \textlatin{layer};}

% TODO: ...

\paragraph{Τι διαφορά θα είχε αν είχαμε 2 ή περισσότερους γραμμικούς μετασχηματισμούς στη σειρά;}

% TODO: ...

\phantomsection
\subsection*{Ζητούμενο 6 -- \textlatin{Forward Pass}}
\addcontentsline{toc}{subsection}{Ζητούμενο 6 -- Forward Pass}

% Συμπληρώστε τα κενά στις θέσεις models.py:EX6 στη μέθοδο __forward__
% Απαντήστε στα παρακάτω ερωτήματα:

\paragraph{Διαισθητική ερμηνεία της αναπαράστασης κέντρο-βάρους:}

% TODO: ...

\paragraph{Πιθανές αδυναμίες της προσέγγισης για αναπαράσταση κειμένων:}

% TODO: ...

\phantomsection
\subsection*{Ζητούμενο 7 -- \textlatin{DataLoaders}}
\addcontentsline{toc}{subsection}{Ζητούμενο 7 -- DataLoaders}

% Συμπληρώστε τα κενά στις θέσεις main.py:EX7
% Απαντήστε στα παρακάτω ερωτήματα:

\paragraph{Συνέπειες μικρών και μεγάλων \textlatin{mini-batches} στην εκπαίδευση:}

% TODO: ...

\paragraph{Γιατί ανακατεύουμε τη σειρά των \textlatin{mini-batches} σε κάθε εποχή;}

% TODO: ...

\phantomsection
\subsection*{Ζητούμενο 8 -- Βελτιστοποίηση}
\addcontentsline{toc}{subsection}{Ζητούμενο 8 -- Βελτιστοποίηση}

% Συμπληρώστε τα κενά στις θέσεις main.py:EX8
% Αναφέρετε τις επιλογές σας (κριτήριο, παράμετροι, optimizer) και αιτιολογήστε.

\paragraph{Κριτήριο (\textlatin{Loss}):} % TODO: BCEWithLogitsLoss / CrossEntropyLoss -- γιατί;

\paragraph{\textlatin{Optimizer}:} % TODO: Adam / Adagrad / RMSProp -- γιατί;

\paragraph{Υπερπαράμετροι:}
\begin{itemize}
  \item \textlatin{Learning rate}: % TODO
  \item \textlatin{Batch size}: % TODO
  \item Εποχές: % TODO
  \item Άλλες: % TODO
\end{itemize}

\phantomsection
\subsection*{Ζητούμενο 9 -- Εκπαίδευση}
\addcontentsline{toc}{subsection}{Ζητούμενο 9 -- Εκπαίδευση}

% Συμπληρώστε τα κενά στις θέσεις training.py:EX9.

\begin{lstlisting}[caption={Σύνοψη εκπαίδευσης}]
% TODO: Επικολλήστε σύνοψη εκπαίδευσης (π.χ. loss ανά εποχή)
\end{lstlisting}

\phantomsection
\subsection*{Ζητούμενο 10 -- Αξιολόγηση \textlatin{Baseline} Μοντέλου}
\addcontentsline{toc}{subsection}{Ζητούμενο 10 -- Αξιολόγηση Baseline Μοντέλου}

% Για κάθε ένα από τα 2 datasets, αναφέρετε accuracy, F1-score (macro), recall (macro).
% Δημιουργήστε γραφικές παραστάσεις training/test loss ανά εποχή.

\begin{table}[H]
  \centering
  \caption{Αποτελέσματα \textlatin{Baseline} μοντέλου}
  \begin{tabular}{lccc}
    \toprule
    \textbf{\textlatin{Dataset}} & \textbf{\textlatin{Accuracy}} & \textbf{\textlatin{F1 (macro)}} & \textbf{\textlatin{Recall (macro)}} \\
    \midrule
    \textlatin{Sentence Polarity} & -- & -- & -- \\
    \textlatin{Semeval 2017}      & -- & -- & -- \\
    \bottomrule
  \end{tabular}
\end{table}

\begin{figure}[H]
  \centering
  % \includegraphics[width=0.8\textwidth]{figures/loss_curve_baseline.png}
  \caption{Καμπύλες εκπαίδευσης (\textlatin{training \& test loss}) ανά εποχή -- \textlatin{Baseline}}
  \label{fig:loss_baseline}
\end{figure}

\phantomsection
\subsection*{Ζητούμενο 11 -- Κατηγοριοποίηση με Χρήση \textlatin{LLM}}
\addcontentsline{toc}{subsection}{Ζητούμενο 11 -- Κατηγοριοποίηση με Χρήση LLM}

% Επιλέξτε τουλάχιστον 20 κείμενα από κάθε κατηγορία και χρησιμοποιήστε το ChatGPT.

\paragraph{\textlatin{Prompts} που χρησιμοποιήθηκαν:}

% TODO: Περιγράψτε τα prompts που δοκιμάσατε

\paragraph{Επίδοση μοντέλου:}

\begin{table}[H]
  \centering
  \caption{Αποτελέσματα \textlatin{ChatGPT}}
  \begin{tabular}{lccc}
    \toprule
    \textbf{\textlatin{Dataset}} & \textbf{\textlatin{Accuracy}} & \textbf{\textlatin{F1 (macro)}} & \textbf{\textlatin{Recall (macro)}} \\
    \midrule
    \textlatin{Sentence Polarity} & -- & -- & -- \\
    \textlatin{Semeval 2017}      & -- & -- & -- \\
    \bottomrule
  \end{tabular}
\end{table}

\paragraph{Ανάλυση αποτελεσμάτων:}

% TODO: Έκανε κάποιο λάθος; Ποιες λέξεις ήταν σημαντικότερες για το συναίσθημα;

% ============================================================
%  ΜΕΡΟΣ Β: ΚΥΡΙΑ ΑΣΚΗΣΗ
% ============================================================
\newpage
\phantomsection
\section*{Μέρος Β΄: Κύρια Άσκηση}
\addcontentsline{toc}{section}{Μέρος Β΄: Κύρια Άσκηση}

Για τα ερωτήματα 1--5 χρησιμοποιήσαμε το \textbf{\textlatin{[Sentence Polarity Dataset / Semeval 2017 Task4-A]}}.
% TODO: επιλέξτε ένα dataset

\phantomsection
\subsection*{Ερώτημα 1 -- \textlatin{Mean + Max Pooling}}
\addcontentsline{toc}{subsection}{Ερώτημα 1 -- Mean + Max Pooling}

\paragraph{1.1 Υλοποίηση:}

% Υπολογισμός αναπαράστασης u = [mean(E) || max(E)]

\paragraph{1.2 Διαφορές από την αρχική αναπαράσταση:}

% TODO: Τι παραπάνω πληροφορία αποτυπώνει η συνένωση mean + max pooling;

\begin{table}[H]
  \centering
  \caption{Αποτελέσματα Ερωτήματος 1 (\textlatin{Mean+Max Pooling})}
  \begin{tabular}{lccc}
    \toprule
    \textbf{Μοντέλο} & \textbf{\textlatin{Accuracy}} & \textbf{\textlatin{F1 (macro)}} & \textbf{\textlatin{Recall (macro)}} \\
    \midrule
    \textlatin{Mean+Max Pooling} & -- & -- & -- \\
    \bottomrule
  \end{tabular}
\end{table}

\phantomsection
\subsection*{Ερώτημα 2 -- \textlatin{LSTM}}
\addcontentsline{toc}{subsection}{Ερώτημα 2 -- LSTM}

\paragraph{2.1 \textlatin{Validation set \& Early Stopping}:}

% TODO: Περιγράψτε τη χρήση EarlyStopper (early stopping μετά από 5 εποχές αύξησης validation loss)

\paragraph{2.2 Χρήση τελευταίας εξόδου \textlatin{LSTM} ($h_N$):}

% TODO: Αποτελέσματα με τη χρήση h_N (χωρίς zero-padded timesteps)

\paragraph{2.3 Αμφίδρομο \textlatin{LSTM (BiLSTM)}:}

% TODO: Αποτελέσματα με bidirectional=True

\begin{table}[H]
  \centering
  \caption{Αποτελέσματα Ερωτήματος 2 (\textlatin{LSTM})}
  \begin{tabular}{lccc}
    \toprule
    \textbf{Μοντέλο} & \textbf{\textlatin{Accuracy}} & \textbf{\textlatin{F1 (macro)}} & \textbf{\textlatin{Recall (macro)}} \\
    \midrule
    \textlatin{LSTM (unidirectional)}  & -- & -- & -- \\
    \textlatin{BiLSTM}                 & -- & -- & -- \\
    \bottomrule
  \end{tabular}
\end{table}

\begin{figure}[H]
  \centering
  % \includegraphics[width=0.8\textwidth]{figures/loss_curve_lstm.png}
  \caption{Καμπύλες εκπαίδευσης \textlatin{LSTM}}
  \label{fig:loss_lstm}
\end{figure}

\phantomsection
\subsection*{Ερώτημα 3 -- \textlatin{Self-Attention}}
\addcontentsline{toc}{subsection}{Ερώτημα 3 -- Self-Attention}

\paragraph{3.1 Υλοποίηση \& Αποτελέσματα (\textlatin{SimpleSelfAttentionModel}):}

% TODO: Συμπληρώστε τα κενά στο attention.SimpleSelfAttentionModel
% Εφαρμόστε average-pooling πριν το τελευταίο layer

\begin{table}[H]
  \centering
  \caption{Αποτελέσματα Ερωτήματος 3 (\textlatin{Self-Attention})}
  \begin{tabular}{lccc}
    \toprule
    \textbf{Μοντέλο} & \textbf{\textlatin{Accuracy}} & \textbf{\textlatin{F1 (macro)}} & \textbf{\textlatin{Recall (macro)}} \\
    \midrule
    \textlatin{Self-Attention} & -- & -- & -- \\
    \bottomrule
  \end{tabular}
\end{table}

\paragraph{3.2 \textlatin{Queries, Keys, Values} και \textlatin{Position Embeddings}:}

% TODO: Τι είναι τα queries, keys, values στην κλάση attention.Head;
% Τι είναι τα position_embeddings στο SimpleSelfAttentionModel;

\phantomsection
\subsection*{Ερώτημα 4 -- \textlatin{Multi-Head Attention}}
\addcontentsline{toc}{subsection}{Ερώτημα 4 -- Multi-Head Attention}

% Συμπληρώστε τα κενά στο attention.MultiHeadAttentionModel
% Βασιστείτε στο SimpleSelfAttentionModel

\begin{table}[H]
  \centering
  \caption{Αποτελέσματα Ερωτήματος 4 (\textlatin{Multi-Head Attention})}
  \begin{tabular}{lccc}
    \toprule
    \textbf{Μοντέλο} & \textbf{\textlatin{Accuracy}} & \textbf{\textlatin{F1 (macro)}} & \textbf{\textlatin{Recall (macro)}} \\
    \midrule
    \textlatin{MultiHead-Attention} & -- & -- & -- \\
    \bottomrule
  \end{tabular}
\end{table}

\phantomsection
\subsection*{Ερώτημα 5 -- \textlatin{Transformer-Encoder}}
\addcontentsline{toc}{subsection}{Ερώτημα 5 -- Transformer-Encoder}

% Συμπληρώστε τα κενά στο attention.TransformerEncoderModel
% Average-pooling για εξαγωγή αναπαράστασης πρότασης

\paragraph{Διαφορά από \textlatin{MultiHeadAttentionModel}:}

% TODO: ...

\paragraph{\textlatin{Default} τιμές παραμέτρων στην κλασική αρχιτεκτονική \textlatin{Transformer}:}

% TODO: ...

\begin{table}[H]
  \centering
  \caption{Αποτελέσματα Ερωτήματος 5 (\textlatin{Transformer-Encoder})}
  \begin{tabular}{lccc}
    \toprule
    \textbf{Μοντέλο} & \textbf{\textlatin{Accuracy}} & \textbf{\textlatin{F1 (macro)}} & \textbf{\textlatin{Recall (macro)}} \\
    \midrule
    \textlatin{Transformer-Encoder} & -- & -- & -- \\
    \bottomrule
  \end{tabular}
\end{table}

\phantomsection
\subsection*{Ερώτημα 6 -- \textlatin{Pre-Trained Transformers (Zero-Shot)}}
\addcontentsline{toc}{subsection}{Ερώτημα 6 -- Pre-Trained Transformers}

% Εκτελέστε το transfer_pretrained.py για τουλάχιστον 3 μοντέλα (από HuggingFace)
% για κάθε dataset. Συγκρίνετε τις επιδόσεις.

\begin{table}[H]
  \centering
  \caption{Αποτελέσματα \textlatin{Pre-Trained Transformers} -- \textlatin{Sentence Polarity}}
  \begin{tabular}{lccc}
    \toprule
    \textbf{Μοντέλο} & \textbf{\textlatin{Accuracy}} & \textbf{\textlatin{F1 (macro)}} & \textbf{\textlatin{Recall (macro)}} \\
    \midrule
    % TODO: \textlatin{Μοντέλο 1 (π.χ. bert-base-uncased)} & -- & -- & -- \\
    % TODO: \textlatin{Μοντέλο 2}                          & -- & -- & -- \\
    % TODO: \textlatin{Μοντέλο 3}                          & -- & -- & -- \\
    \bottomrule
  \end{tabular}
\end{table}

\begin{table}[H]
  \centering
  \caption{Αποτελέσματα \textlatin{Pre-Trained Transformers} -- \textlatin{Semeval 2017}}
  \begin{tabular}{lccc}
    \toprule
    \textbf{Μοντέλο} & \textbf{\textlatin{Accuracy}} & \textbf{\textlatin{F1 (macro)}} & \textbf{\textlatin{Recall (macro)}} \\
    \midrule
    % TODO: \textlatin{Μοντέλο 1} & -- & -- & -- \\
    % TODO: \textlatin{Μοντέλο 2} & -- & -- & -- \\
    % TODO: \textlatin{Μοντέλο 3} & -- & -- & -- \\
    \bottomrule
  \end{tabular}
\end{table}

\paragraph{Σχολιασμός αποτελεσμάτων:}

% TODO: Συγκρίνετε τα μοντέλα μεταξύ τους και με τα προηγούμενα αποτελέσματα.

\phantomsection
\subsection*{Ερώτημα 7 -- \textlatin{Fine-Tuning Pre-Trained Transformers}}
\addcontentsline{toc}{subsection}{Ερώτημα 7 -- Fine-Tuning Pre-Trained Transformers}

% Χρησιμοποιήστε finetune_pretrained.py (Google Colab)
% Τουλάχιστον 3 μοντέλα για κάθε dataset

\begin{table}[H]
  \centering
  \caption{Αποτελέσματα \textlatin{Fine-Tuned Transformers} -- \textlatin{Sentence Polarity}}
  \begin{tabular}{lccc}
    \toprule
    \textbf{Μοντέλο} & \textbf{\textlatin{Accuracy}} & \textbf{\textlatin{F1 (macro)}} & \textbf{\textlatin{Recall (macro)}} \\
    \midrule
    % TODO: \textlatin{Μοντέλο 1} & -- & -- & -- \\
    % TODO: \textlatin{Μοντέλο 2} & -- & -- & -- \\
    % TODO: \textlatin{Μοντέλο 3} & -- & -- & -- \\
    \bottomrule
  \end{tabular}
\end{table}

\begin{table}[H]
  \centering
  \caption{Αποτελέσματα \textlatin{Fine-Tuned Transformers} -- \textlatin{Semeval 2017}}
  \begin{tabular}{lccc}
    \toprule
    \textbf{Μοντέλο} & \textbf{\textlatin{Accuracy}} & \textbf{\textlatin{F1 (macro)}} & \textbf{\textlatin{Recall (macro)}} \\
    \midrule
    % TODO: \textlatin{Μοντέλο 1} & -- & -- & -- \\
    % TODO: \textlatin{Μοντέλο 2} & -- & -- & -- \\
    % TODO: \textlatin{Μοντέλο 3} & -- & -- & -- \\
    \bottomrule
  \end{tabular}
\end{table}

\paragraph{Σχολιασμός αποτελεσμάτων:}

% TODO: Συγκρίνετε zero-shot vs fine-tuned. Ποια μοντέλα απέδωσαν καλύτερα και γιατί;

\phantomsection
\subsection*{Ερώτημα 8 -- Διάλογος με \textlatin{ChatGPT} για \textlatin{GPT} Κώδικα}
\addcontentsline{toc}{subsection}{Ερώτημα 8 -- Διάλογος με ChatGPT}

% Χρησιμοποιήστε τον κώδικα "Let's build GPT: from scratch, in code, spelled out"
% Ζητήστε από το ChatGPT: εξήγηση, αξιολόγηση, refactoring

\paragraph{Παραδείγματα διαλόγου:}

% TODO: Επικολλήστε αντιπροσωπευτικά αποσπάσματα του διαλόγου με το ChatGPT

\paragraph{Σχολιασμός απόδοσης:}

% TODO: Εξηγήθηκε σωστά ο κώδικας; Ήταν η αξιολόγηση σωστή; Πώς ήταν το refactoring;

% ============================================================
%  ΣΥΓΚΕΝΤΡΩΤΙΚΑ ΑΠΟΤΕΛΕΣΜΑΤΑ
% ============================================================
\newpage
\phantomsection
\section*{Συγκεντρωτικά Αποτελέσματα}
\addcontentsline{toc}{section}{Συγκεντρωτικά Αποτελέσματα}

\begin{table}[H]
  \centering
  \caption{Σύγκριση όλων των μοντέλων -- \textbf{\textlatin{[Επιλεγμένο Dataset]}}}
  \begin{tabular}{lccc}
    \toprule
    \textbf{Μοντέλο} & \textbf{\textlatin{Accuracy}} & \textbf{\textlatin{F1 (macro)}} & \textbf{\textlatin{Recall (macro)}} \\
    \midrule
    \textlatin{Baseline DNN (mean pooling)}  & -- & -- & -- \\
    \textlatin{Mean + Max Pooling}           & -- & -- & -- \\
    \textlatin{LSTM}                         & -- & -- & -- \\
    \textlatin{BiLSTM}                       & -- & -- & -- \\
    \textlatin{Self-Attention}               & -- & -- & -- \\
    \textlatin{Multi-Head Attention}         & -- & -- & -- \\
    \textlatin{Transformer-Encoder}          & -- & -- & -- \\
    \textlatin{Pre-Trained (best)}           & -- & -- & -- \\
    \textlatin{Fine-Tuned (best)}            & -- & -- & -- \\
    \bottomrule
  \end{tabular}
\end{table}

\end{document}
